\subsection{The full level-seven modular-form determinant}
\label{sec:gl2-minus7}

\begin{proposition}[K17-MINUS7-GL2]\label{prop:minus7-gl2}
Let \(\ell_0=25078104767246023\), and let \(T_2\) be the usual
Hecke operator on the entire \(61\)-dimensional space
\(S_{33}(\Gamma_1(7))\). Put \(t_2=1+2^{32}=4294967297\).
On its full \(\ell_0\)-integral Fourier lattice one has
\[
 \det(t_2^3I-T_2^3)
    \equiv15133664954664570\not\equiv0\pmod{\ell_0}.
 \tag{G7.1}
\]
Consequently \(t_2-\zeta a_2(f)\) is a unit at every place above
\(\ell_0\), for every normalized eigenform in this space and every
\(\zeta\in\mu_3\).
\end{proposition}

\begin{proof}
We first identify the entire space and its integral lattice.
Fix \(\xi=\exp(2\pi i/6)\) and the mod-\(7\) character
\(\chi(3)=\xi\). Extend every such character by zero on multiples
of \(7\), and write \(Q=\exp(2\pi i\tau)\). Define
\[
 E_r=-\frac1{14}\sum_{u=1}^6u\chi(u)^r
     +\sum_{n\ge1}\left(\sum_{d\mid n}\chi(d)^r\right)Q^n,
 \qquad r=1,3,5,
\]
and put \(a=E_1,\ b=E_3,\ c=E_5\).
These are exactly the weight-one Eisenstein series, with their
diamond characters, in
\cite[\S2, pp.~8--9]{K17MeierOzornova2020}. Their constants are
\[
 a(0)=\frac{2\xi+1}{7},\qquad b(0)=\frac12,\qquad
 c(0)=\frac{3-2\xi}{7}.
 \tag{G7.2}
\]
The integral conic-ring theorem
\cite[Proposition~2.4 and its proof, pp.~10--11]{K17MeierOzornova2020}
gives
\[
 M_*(\Gamma_1(7);\mathbf Z[1/7])
  =\mathbf Z[1/7][z_1,z_2,z_3]/
       (z_1z_2+z_2z_3+z_3z_1).
 \tag{G7.3}
\]
Here the \(z_i\) are the explicit integral weight-one basis of
their Lemma~2.3. Their basis-change matrix from \((a,b,c)\)
has determinant
\[
 \frac2{27}(84\xi-42)=\frac{28}{9}(2\xi-1).
 \tag{G7.4}
\]
This is a unit after inverting \(2,3,7\), since
\((2\xi-1)^2=-3\).
The unique quadratic relation is diamond-invariant
\cite[Lemma~2.6, p.~11]{K17MeierOzornova2020}.
The invariant quadratic monomials in \(a,b,c\) are \(ac\) and
\(b^2\); comparing (G7.2) gives
\[
 ac=\frac47b^2.
 \tag{G7.5}
\]
Thus, over \(R=\mathbf Z[1/42,\xi]\), the full ring is
\(R[a,b,c]/(ac-\tfrac47b^2)\).
This is also the ring after localization at either prime above
\(\ell_0\) and after completion: the flat base-change statement
is \cite[Lemma~2.1, p.~7]{K17MeierOzornova2020}.
The monic relation (G7.5) gives the following \(R\)-basis in
weight \(30\):
\[
 \{a^{30-j}b^j:0\le j\le30\}
 \ \cup\
 \{b^{30-j}c^j:1\le j\le30\}.
 \tag{G7.6}
\]
For example, division by the single relation with leading
monomial \(ac\) leaves exactly these standard monomials, so
there is no unstated saturation or finite-index assertion.

Now put
\[
 D(\tau)=\eta(\tau)^3\eta(7\tau)^3
        =Q\prod_{n\ge1}(1-Q^n)^3(1-Q^{7n})^3.
 \tag{G7.7}
\]
The eta-quotient theorem
\cite[Theorem~1.1 and Remark~1.2, p.~2]{K17SavittNewman}
applies with exponents \(r_1=r_7=3\). The two required sums
are both \(24\), its weight is \(3\), and its nebentypus is
\(\bigl(\frac{-7^3}{\,\cdot\,}\bigr)=\chi_{-7}=\chi^3\).
It has no zeros in the upper half-plane, by its convergent
product. It has order one at the infinity cusp of
\(\Gamma_0(7)\). The explicit eta transformation
\cite[Proposition~3.1 and its proof, p.~4]{K17SavittNewman}
gives, with determinant-normalized weight-three slash,
\[
 D|_3\begin{pmatrix}0&-1\\7&0\end{pmatrix}=iD.
 \tag{G7.8}
\]
It therefore has order one at the other cusp \(0\), too.

For completeness, both orders remain one after passage to
\(\Gamma_1(7)\). The projective quotient of
\(\Gamma_0(7)\) by \(\Gamma_1(7)\) has order three.
At infinity its width is one in both groups, since \(T\)
already lies in \(\Gamma_1(7)\); at zero its width is seven
in both groups, since \(ST^7S^{-1}\) already lies there.
Thus each of the two \(\Gamma_0(7)\) cusps has three
preimages with ramification index one. These are all six
cusps, and \(D\) vanishes simply at every one of them.
Consequently division by \(D\) is an isomorphism
\[
 S_{33}(\Gamma_1(7))=D\,M_{30}(\Gamma_1(7)).
 \tag{G7.9}
\]
It is an equality of the full integral Fourier lattices
at \(\ell_0\) as well. Indeed \(D/Q\) and its inverse both
belong to \(1+Q\mathbf Z[[Q]]\). If a cusp form has integral
Fourier coefficients, its holomorphic quotient by \(D\)
therefore has integral Fourier coefficients. By (G7.3)--(G7.6)
that quotient belongs to the full displayed modular lattice.
Conversely multiplication by \(D\) preserves integral
coefficients. This proves that \(D\) times (G7.6) is an
actual integral basis, of size \(61\).

Both families of this basis have nebentypus
\(\chi^{3+2j}\), where \(j\) is their displayed index.
Their three odd-character dimensions are \(20,21,20\),
for \(\chi,\chi^3,\chi^5\), respectively.
On the \(\chi^r\) summand the exact Hecke coefficient formula is
\[
 [Q^n]T_2F=[Q^{2n}]F+
       \chi(2)^r\,2^{32}\,
       \begin{cases}[Q^{n/2}]F,&2\mid n,\\0,&2\nmid n.\end{cases}
 \tag{G7.10}
\]
It shows directly that the full integral Fourier lattice is
\(T_2\)-stable.

We describe the complete finite calculation of (G7.1).
The standalone program
\texttt{hecke\_minus7/gl2/gl2\_minus7\_residue\_certificate.py}
reconstructs all \(61\) series from (G7.2), the divisor sums,
and (G7.7), to precision \(Q^{161}\).
It works in \(\mathbf F_{\ell_0}\), using
\(\xi=3287403608222320\), and verifies
\(\xi^2-\xi+1=0\), \(\xi^3=-1\), and the conic relation.
The prime \(\ell_0\) is certified by the separate exact
recursive Lucas certificate, with its complete factorization
of \(\ell_0-1\); no probable-prime assertion is used.
For each character block, the first \(d\) positive
Fourier coefficients give a nonsingular identifying matrix
of size \(d\). The following are the exact residues:
\[
\begin{array}{c|r|r|r}
 r&d&\det(\text{identifying matrix})
       &\det(t_2^3I-T_2^3)\bmod\ell_0\\ \hline
 1&20&13431297379939036&10804607418746181\\
 3&21& 7862163528894746& 2808906210861363\\
 5&20&21251185086201699&10804607418746181
\end{array}
 \tag{G7.11}
\]
The program solves for \(T_2\) using these unit minors and
(G7.10), then checks every Fourier coefficient through
\(Q^{80}\). The finite matrix is the reduction of the actual
operator because the basis is full, the operator preserves
its lattice, and the identifying coefficient map is
injective modulo \(\ell_0\). The additional coefficient
checks are independent consistency checks; they are not
being used in place of this identification.
The determinant is evaluated by exact finite-field
elimination. The program separately verifies on each
block the identity
\[
 \det(t_2^3I-T_2^3)
  =\prod_{j=0}^2\det(t_2I-\xi^{2j}T_2).
 \tag{G7.12}
\]
Multiplying the three last-column residues in (G7.11)
gives \(15133664954664570\), proving (G7.1).
The full matrices, identifying minors, raw Hecke
coefficients, and both determinant calculations are
recorded in the accompanying exact JSON certificate.
An independent implementation,
\texttt{gl2\_minus7\_integral\_archive\_certificate.py},
uses the integral \(z_1,z_2,z_3\) basis and the cyclic
diamond action on the full \(61\)-dimensional space, with
no sixth root of unity or character decomposition. Its
first \(61\) Fourier coefficients give identifying
determinant \(-1\), and fraction-free elimination gives
the same residue \(15133664954664570\).

There are two choices of \(\xi\) at the two primes above
\(\ell_0\), since \(\ell_0\equiv1\pmod6\).
Conjugation exchanges the \(r=1\) and \(r=5\) blocks
and permutes the three twists. The total determinant is
unchanged. Equivalently it is the determinant on the full
rational space with all odd characters, so the displayed
unit conclusion holds for either choice.
Finally all Hecke eigenvalues are integral on the lattice
just proved. The product (G7.12) being a unit makes every
individual factor \(t_2-\zeta a_2(f)\) a unit after any
coefficient extension. This proves the last assertion.
\end{proof}
